\subsection{DEFINITION OF PROJECTIVE SPACES}

DEFINITION 4. Given a left (resp. right) vector space V over a field K, the left (resp. right) projective space derived from V, denoted by P(V), is the quotient of the complement V - \{0\} of \{0\} in V by the equivalence relation Δ(V) ``there exists λ ≠ 0 in K such that y = λx (resp. y = xλ) between x and y in V - \{0\}.

When \(\mathbf{V} = \mathrm{K}_s^{n + 1}\) , we also write \(\mathbf{P}_n(\mathbf{K})\) instead of \(\mathbf{P}(\mathrm{K}_s^{n + 1})\) and \(\Delta_n(\mathbf{K})\) instead of \(\Delta (\mathbf{V})\) .

Definition 4 can also be expressed by saying that \(\mathbf{P}(\mathbf{V})\) is the set of lines (passing through 0) in V with the origin removed; \(\mathbf{P}(\mathbf{V})\) is therefore canonically identified with the set of lines (passing through 0) in V. The elements of a projective space are called the points of that space.

When V is of dimension n, the integer n - 1 is called the dimension of the projective space \(\mathbf{P}(\mathbf{V})\) if n is finite, and the cardinal n otherwise; this cardinal is denoted by \(\dim_{\mathbb{K}} \mathbf{P}(\mathbf{V})\) or \(\dim \mathbf{P}(\mathbf{V})\) . Thus a projective space of dimension -1 is empty and a projective space of dimension 0 is a single point. A projective space of dimension 1 (resp. 2) is called a projective line (resp. projective plane).

Henceforth we shall only consider left projective spaces.

